% ================= APPENDIX: FINAL SPECIFICATION SHEET =================
% \appendix is declared in main.tex, before the first appendix input.
\section{Final Specification Sheet}
\label{app:specs}

This appendix is the single place where the finished design is stated as a set
of numbers: what the cube is, what it weighs, what the wheel is, and what
operating point it was designed to. Every figure below is \emph{quoted}, never
derived --- the last column of each table names where it comes from, and the
derivation stays in the section that owns it. Nothing here is new work, and
nothing here may disagree with its source: all recurring values are drawn from
the single definition file \texttt{sections/00\_values.tex}, which in turn
mirrors the run of \texttt{analysis/SIZING.py}, so a value cannot be updated in
one table and left stale in another.

Two distinct masses appear throughout and are not interchangeable. The
\emph{sizing basis}, \qty{\Mbasis}{\kilo\gram}, is the mass the reaction wheels
were designed against and the mass at which PER-06 is derived; it is the
heaviest case the design admits. The \emph{as-built} mass,
\qty{\Mbuilt}{\kilo\gram}, is what the integrated article weighs. The
\qty{\Mdelta}{\gram} between them is realised margin, quantified in
Appendix~\ref{app:cad_mass}, and is deliberately not folded back into the
sizing.

\begin{figure}[htbp]
\centering
\adjincludegraphics[trim={0.34\width} {0.06\height} {0.34\width} {0.06\height},
                    clip, width=0.38\textwidth]{images/armonia_assembly_s1.jpg}
\caption[Configuration of the integrated cube]{Configuration the tables below
describe: an outer frame of edge rails and corner contact features enclosing
the inner structure, the three orthogonal reaction-wheel modules and the
electronics. The dimensioned general-assembly drawing is
Figure~\ref{fig:dwg_3d}.}
\label{fig:spec_config}
\end{figure}

\subsection{Configuration and Principal Dimensions}
\label{app:specs_dims}

\begin{table}[htbp]
\centering
\caption[Principal dimensions and configuration]{Principal dimensions and
configuration of the three-axis cube.}
\label{tab:spec_dims}
\small
\setlength{\tabcolsep}{4pt}
\begin{tabular}{@{} l L{5.4cm} c L{3.0cm} @{}}
\toprule
Symbol & Quantity & Value & Source \\
\midrule
\multicolumn{4}{@{}l}{\textbf{Body}}\\
\midrule
$L$        & Outer edge length (frozen)            & \qty{\cubeL}{\milli\metre}\textsuperscript{*} & CON-01, gate M2a \\
$t_w$      & Frame wall thickness                  & \qty{\cubeWall}{\milli\metre} & Tab.~\ref{tab:cad_params} \\
$M$        & Total mass, sizing basis              & \qty{\Mbasis}{\kilo\gram} & Tab.~\ref{tab:jumpup_budget} \\
$M$        & Total mass, as built                  & \qty{\Mbuilt}{\kilo\gram} & Tab.~\ref{tab:mass_verification} \\
           & Architecture                          & Two-part: outer frame $+$ inner structure & Sec.~\ref{sec:mechanical} \\
           & Corner contact features               & 8 & Tab.~\ref{tab:cad_parts} \\
\midrule
\multicolumn{4}{@{}l}{\textbf{Reaction wheel} ($\times3$, mutually orthogonal)}\\
\midrule
$D_w$      & Ring outer diameter                   & \qty{\wheelOD}{\milli\metre} & Tab.~\ref{tab:rw_params} \\
$D_i$      & Ring inner diameter                   & \qty{\wheelID}{\milli\metre} & Tab.~\ref{tab:rw_params} \\
$w_r$      & Ring radial width (= hole depth)      & \qty{\wheelWr}{\milli\metre} & Tab.~\ref{tab:rw_params} \\
$t$        & Wheel axial thickness                 & \qty{\wheelT}{\milli\metre} & Tab.~\ref{tab:rw_params} \\
$n_s, w_s$ & Spoke count and width                 & $\wheelSpokes$, \qty{\wheelSpokeW}{\milli\metre} & Tab.~\ref{tab:rw_params} \\
$N$        & Ballast stations populated (of $\wheelNmax$) & \num{\wheelN} & Tab.~\ref{tab:rw_params} \\
$d_h$      & Ballast hole diameter (M6 clearance)  & \qty{\wheelHole}{\milli\metre} & Tab.~\ref{tab:rw_params} \\
$R_{\text{mean}}$ & Ballast centroid radius        & \qty{\wheelRmean}{\milli\metre} & Tab.~\ref{tab:rw_params} \\
\midrule
\multicolumn{4}{@{}l}{\textbf{Actuation and power}}\\
\midrule
           & Motor ($\times3$)                     & T-Motor MN4006 KV380 & Tab.~\ref{tab:bom-supplied} \\
           & Driver ($\times3$)                    & mjbots moteus-n1 & Tab.~\ref{tab:bom-supplied} \\
           & Brake servo ($\times3$)               & TowerPro MG92B & Tab.~\ref{tab:bom-supplied} \\
           & Battery                               & 6S LiPo, \qty{22.2}{\volt}, \qty{1550}{\milli\ampere\hour} & Tab.~\ref{tab:bom-supplied} \\
           & Flight computer                       & Teensy 4.1 & Tab.~\ref{tab:bom-supplied} \\
\bottomrule
\end{tabular}

\vspace{0.4em}
{\footnotesize\RaggedRight\textsuperscript{*}\,The sizing analysis of
Section~\ref{sec:sizing} is run at \qty{\cubeLan}{\milli\metre}, one millimetre
under the frozen value. The discrepancy is material at the present design
point: since $h_w \propto ML^{3/2}$, closing at \qty{\cubeL}{\milli\metre}
raises $I_{w,\text{target}}$ by \qty{1.22}{\percent} and leaves the selected
wheel at \qty{99.3}{\percent} of its own requirement rather than
\qty{\marginSized}{\percent}. The reconciliation is an open decision, recorded
in Section~\ref{sec:requirements}, and is to be taken before the wheel is
committed to manufacture.\par}
\end{table}

\subsection{Mass Decomposition}
\label{app:specs_mass}

Table~\ref{tab:spec_mass} is the itemised roll-up behind the six-line summary
of Table~\ref{tab:massbudget}. It is grouped by the same categories, so the two
reconcile line for line. The \emph{basis} column is the important one: it
separates what the design knows (weighed parts and datasheet figures) from what
it still projects (the structural allowance), which is \qty{24.0}{\percent} of
the projected total and the dominant uncertainty in the budget.

\begin{small}
\setlength{\tabcolsep}{4pt}
\begin{longtable}{@{} L{6.2cm} r c r L{2.6cm} @{}}
\caption[Itemised mass decomposition]{Itemised mass decomposition at the
converged design point, rolled up by \texttt{analysis/SIZING.py} (Stage~6).
Quantities are per assembled cube.}
\label{tab:spec_mass}\\
\toprule
Item & Unit (g) & Qty & Total (g) & Basis \\
\midrule
\endfirsthead
\multicolumn{5}{c}{\footnotesize\itshape \tablename~\thetable\ (continued)}\\
\toprule
Item & Unit (g) & Qty & Total (g) & Basis \\
\midrule
\endhead
\midrule
\multicolumn{5}{r}{\footnotesize\itshape Continued on next page}\\
\endfoot
\bottomrule
\endlastfoot

\multicolumn{5}{@{}l}{\textbf{Wheels} --- \qty{702.4}{\gram}, \qty{39.0}{\percent}}\\
\midrule
Reaction wheel (PET-CF, \qty{\wheelOD}{\milli\metre}, $\wheelN\times$ M6 radial) & \wheelMass & 3 & 702.42 & Sized, Sec.~\ref{sec:rw_sizing} \\
\midrule

\multicolumn{5}{@{}l}{\textbf{Motors and brake servos} --- \qty{245.4}{\gram}, \qty{13.6}{\percent}}\\
\midrule
T-Motor MN4006 KV380        & 68.00 & 3 & 204.00 & Datasheet \\
TowerPro MG92B brake servo  & 13.80 & 3 &  41.40 & Datasheet \\
\midrule

\multicolumn{5}{@{}l}{\textbf{Power} --- \qty{279.0}{\gram}, \qty{15.5}{\percent}}\\
\midrule
Tattu R-Line V5.0 6S \qty{1550}{\milli\ampere\hour} LiPo & 254.00 & 1 & 254.00 & As weighed \\
XT60 connector pair         & 15.00 & 1 &  15.00 & As weighed \\
LM2596 step-down regulator  & 10.00 & 1 &  10.00 & As weighed \\
\midrule

\multicolumn{5}{@{}l}{\textbf{Electronics} --- \qty{95.2}{\gram}, \qty{5.3}{\percent}}\\
\midrule
mjbots moteus-n1 driver     & 14.60 & 3 &  43.80 & Datasheet \\
mjbots MA600 breakout $+$ magnet & 6.00 & 3 & 18.00 & Datasheet \\
Teensy 4.1                  & 11.00 & 1 &  11.00 & Datasheet \\
CAN-FD adapter for Teensy 4.1 &  6.00 & 1 &  6.00 & Datasheet \\
mjcanfd-usb-1x              &  3.40 & 1 &   3.40 & Datasheet \\
Seeed XIAO ESP32-C6 $+$ antenna & 8.00 & 1 & 8.00 & Datasheet \\
SparkFun BMI270 IMU (Qwiic) &  5.00 & 1 &   5.00 & Datasheet \\
\midrule

\multicolumn{5}{@{}l}{\textbf{Wiring, perfboard and passives} --- \qty{45.5}{\gram}, \qty{2.5}{\percent}}\\
\midrule
Protoboard \qtyproduct{15 x 9}{\centi\metre} double-sided & 25.00 & 1 & 25.00 & As weighed \\
JST PH3 cable               &  3.00 & 3 &   9.00 & Estimated \\
100\,\unit{\micro\farad} \qty{16}{\volt} electrolytic cap & 1.00 & 3 & 3.00 & Estimated \\
100\,\unit{\nano\farad} \qty{50}{\volt} cap & 1.00 & 3 & 3.00 & Estimated \\
1N5819 Schottky diode       &  3.00 & 1 &   3.00 & Estimated \\
4.7\,\unit{\nano\farad} cap (CAN) & 1.00 & 2 & 2.00 & Estimated \\
\qty{60.4}{\ohm} 1/2\,W resistor (CAN) & 0.25 & 2 & 0.51 & Estimated \\
\midrule

\multicolumn{5}{@{}l}{\textbf{Structure} --- \qty{\massStructAllow}{\gram}, \qty{24.0}{\percent}}\\
\midrule
Outer frame / contact features & 212.80 & 1 & 212.80 & \textbf{Allowance} \\
Inner structure / motor subframe & 176.90 & 1 & 176.90 & \textbf{Allowance} \\
Fasteners / wiring harness  & 42.80 & 1 &  42.80 & \textbf{Allowance} \\
\end{longtable}
\end{small}

\begin{table}[htbp]
\centering
\caption[Mass reconciliation]{Mass reconciliation. The budget closes on the
sizing basis by construction; the weighed article is lighter, and the
difference is realised margin rather than an unclosed loop.}
\label{tab:spec_mass_rollup}
\small
\begin{tabular}{@{} L{7.2cm} r l @{}}
\toprule
Line & Mass & Status \\
\midrule
Specified hardware (wheels, motors, power, electronics, wiring) & \qty{\massKnown}{\gram} & Measured or datasheet \\
Structure allowance                          & \qty{\massStructAllow}{\gram} & Pending CAD closure \\
\midrule
\textbf{Budget total} $=$ sizing basis        & \textbf{\qty{\Mbasis}{\kilo\gram}} & Converged, $\delta=\qty{0.0}{\gram}$ \\
\textbf{As built, weighed}                    & \textbf{\qty{\Mbuilt}{\kilo\gram}} & App.~\ref{app:cad_mass} \\
\quad implied structure mass                  & \qty{\massStructBuilt}{\gram} & Outcome, to be confirmed by CAD \\
\quad realised margin                         & $-$\qty{\Mdelta}{\gram} & Not spent \\
\midrule
CON-02 cap                                    & \qty{\Mcap}{\kilo\gram} & Met with \qty{\Mdelta}{\gram} to spare \\
\bottomrule
\end{tabular}
\end{table}

\subsection{Reaction Wheel}
\label{app:specs_wheel}

\begin{table}[htbp]
\centering
\caption[Reaction-wheel specification]{Reaction-wheel specification, per wheel.
The wheel is a printed ring on three spokes, tuned to its inertia target by M6
bolt-and-nut sets in radial ID-to-OD through-holes.}
\label{tab:spec_wheel}
\small
\begin{tabular}{@{} L{6.4cm} r L{3.2cm} @{}}
\toprule
Quantity & Value & Source \\
\midrule
Ring material                        & PET-CF, \qty{1290}{\kilo\gram\per\metre\cubed} & Tab.~\ref{tab:rw_params} \\
Mass, unballasted ring $+$ spokes    & \qty{\wheelMassSolid}{\gram} & \texttt{SIZING.py} \\
\quad $+$ M6 steel, $\wheelN$ stations & \qty{\wheelSteel}{\gram} & \texttt{SIZING.py} \\
\quad $-$ plastic displaced, net PET-CF & \qty{\wheelPlastic}{\gram} & \texttt{SIZING.py} \\
\textbf{Mass as ballasted}           & \textbf{\qty{\wheelMass}{\gram}} & Tab.~\ref{tab:rw_params} \\
Inertia, unballasted                 & \qty{\wheelIsolid}{\kilogram\metre\squared} & \texttt{SIZING.py} \\
Net inertia per station              & \qty{\wheelIstation}{\kilogram\metre\squared} & Tab.~\ref{tab:rw_params} \\
\textbf{Inertia as ballasted} $I_w$  & \textbf{\qty{\wheelI}{\kilogram\metre\squared}} & Tab.~\ref{tab:rw_params} \\
\midrule
Requirement $I_{w,\text{target}}$ at the sizing basis & \qty{\Itarget}{\kilogram\metre\squared} & PER-06 \\
\quad margin delivered               & \qty{\marginSized}{\percent} & Tab.~\ref{tab:inertia_cases} \\
Requirement at the as-built mass     & \qty{\ItargetBuilt}{\kilogram\metre\squared} & App.~\ref{app:cad_mass} \\
\quad margin delivered               & \qty{\marginBuilt}{\percent} & App.~\ref{app:cad_mass} \\
Speed at which the built wheel still closes & \qty{\rpmFloor}{\rpm} & App.~\ref{app:cad_mass} \\
\midrule
Centrifugal load per station at \qty{\rpmDesign}{\rpm} & \qty{\wheelForce}{\newton} & \texttt{SIZING.py} \\
\quad against M6 class 8.8 proof load & \qty{0.57}{\percent} & \texttt{SIZING.py} \\
Edge-to-edge gap between stations    & \qty{14.39}{\milli\metre} & \texttt{SIZING.py} \\
Axial wall, each face                & \qty{6.80}{\milli\metre} & \texttt{SIZING.py} \\
\bottomrule
\end{tabular}
\end{table}

\subsection{Operating Point}
\label{app:specs_ops}

\begin{table}[htbp]
\centering
\caption[Design operating point]{Design operating point and the derived
performance figures that follow from it. Values at the sizing basis unless
noted.}
\label{tab:spec_ops}
\small
\begin{tabular}{@{} L{6.4cm} r L{3.2cm} @{}}
\toprule
Quantity & Value & Source \\
\midrule
Design wheel speed $\omega_{\max}$   & \qty{\rpmDesign}{\rpm} & Tab.~\ref{tab:jumpup_budget} \\
\quad as a fraction of the derived ceiling & \qty{\rpmCeilPc}{\percent} & PER-08 \\
Stored momentum per wheel $h_w$      & \qty{\hwDesign}{\kilo\gram\metre\squared\per\second} & Tab.~\ref{tab:jumpup_budget} \\
Stored energy per wheel              & \qty{\wheelKE}{\joule} & Sec.~\ref{sec:jumpup_target} \\
Brake torque $\tau_b$ (sizing value) & \qty{\tauBrake}{\newton\metre} & Tab.~\ref{tab:jumpup_budget} \\
Tip-over torque $\tau_g = MgL/2$     & \qty{\tauGrav}{\newton\metre} & Tab.~\ref{tab:jumpup_budget} \\
\quad at the as-built mass           & \qty{\tauGravBuilt}{\newton\metre} & App.~\ref{app:cad_mass} \\
Brake amplification $\beta$          & \num{\betaBrake} & Tab.~\ref{tab:jumpup_budget} \\
\midrule
Spin-up to design speed at \qty{9}{\ampere} & \qty{\tspin}{\second} & PER-12 \\
\quad floor set by $\tau_g$          & \qty{\tspinFloor}{\second} & Sec.~\ref{sec:jumpup_target} \\
Brake arrest time (requirement)      & $\leq$ \qty{100}{\milli\second} & PER-10 \\
Control and estimation loop rate     & \qty{1}{\kilo\hertz} & PER-05 \\
Driver continuous current            & \qty{9}{\ampere} & PER-09 \\
Motor peak current (short duration)  & \qty{16}{\ampere} & Tab.~\ref{tab:bom-supplied} \\
Motor torque constant $k_t = 60/(2\pi K_v)$ & \qty{\ktMotor}{\newton\metre\per\ampere} & $K_V$, Tab.~\ref{tab:bom-supplied} \\
Drive torque at \qty{9}{\ampere} continuous & \qty{\tauDriveCont}{\newton\metre} & Sec.~\ref{sec:jumpup_target} \\
Pack energy                          & ${\approx}\,\qty{35}{\watt\hour}$ & Tab.~\ref{tab:bom-supplied} \\
\bottomrule
\end{tabular}
\end{table}

\paragraph{What this sheet does not yet carry.} The plant parameters of
Table~\ref{tab:params} --- the body inertia $I_b$ about the pivot, the
centre-of-mass distances $l$ and $l_b$, the damping coefficients $C_b$ and
$C_w$ --- have determination methods assigned but no measured values, and are
recorded here as outstanding rather than estimated. They close on the CAD mass
properties and the 1-DoF identification run of Section~\ref{sec:paramid}, and
belong in this sheet once measured. The same applies to the modelled-versus-
measured wheel mass and inertia rows of Table~\ref{tab:mass_verification},
which are the check that the built wheel is the wheel specified above.
